% Chapter 4, Section 3 _Linear Algebra_ Jim Hefferon
%  http://joshua.smcvt.edu/linearalgebra
%  2001-Jun-12
\section{Laplace 公式}
\textit{本节为选读内容。
  依赖这一材料的后续小节
  只有 Five.III。}

行列式是许多有趣而巧妙的公式的源泉。
下面就是其中之一，它常被用来
手工计算行列式。



\subsectionoptional{Laplace 展开}
\index{Laplace determinant expansion|(}

%<*dis:LaplaceIntro0>
在 $\deter{T}$ 的置换展开中，
每一项都恰从~$T$ 的每一行、每一列中各取一个元素。
\begin{equation*}
  \deter{T}=\!\!
  \sum_{\text{permutations\ }\phi}\!\!\!\!
     t_{1,\phi(1)}t_{2,\phi(2)}\cdots t_{n,\phi(n)}
                                 \deter{P_{\phi}}
\end{equation*}
下面是该公式的 $\nbyn{3}$ 情形。
\begin{align*}
  \begin{vmat}
    t_{1,1}  &t_{1,2}  &t_{1,3} \\
    t_{2,1}  &t_{2,2}  &t_{2,3} \\
    t_{3,1}  &t_{3,2}  &t_{3,3} 
  \end{vmat}
  &=
  \begin{aligned}[t]
    &t_{1,1}t_{2,2}t_{3,3}\deter{P_{\phi_1}}
     +t_{1,1}t_{2,3}t_{3,2}\deter{P_{\phi_2}}
     +t_{1,2}t_{2,1}t_{3,3}\deter{P_{\phi_3}} \\
    &\hspace*{0em}
     +t_{1,2}t_{2,3}t_{3,1}\deter{P_{\phi_4}}
     +t_{1,3}t_{2,1}t_{3,2}\deter{P_{\phi_5}}
     +t_{1,3}t_{2,2}t_{3,1}\deter{P_{\phi_6}}
  \end{aligned}                                      \\
  &=
  \begin{aligned}[t]
    &t_{1,1}t_{2,2}t_{3,3}
     -t_{1,1}t_{2,3}t_{3,2}
     -t_{1,2}t_{2,1}t_{3,3}  \\
    &\hspace*{0em}
     +t_{1,2}t_{2,3}t_{3,1} 
     +t_{1,3}t_{2,1}t_{3,2}
     -t_{1,3}t_{2,2}t_{3,1}
  \end{aligned}
\end{align*}
%</dis:LaplaceIntro0>
%<*dis:LaplaceIntro1>
我们可以聚焦于任意一行或任意一列；这里我们
把来自 $T$ 第一行的元素提取出来。
\begin{align*}
  &=
    t_{1,1}\cdot\big[t_{2,2}t_{3,3}
                      -t_{2,3}t_{3,2}\big]  \\
    &\quad
     -t_{1,2}\cdot\big[t_{2,1}t_{3,3}
                       -t_{2,3}t_{3,1}\big]  \\
    &\quad
     +t_{1,3}\cdot\big[t_{2,1}t_{3,2}
                       -t_{2,2}t_{3,1}\big]
  \tag{$*$}
\end{align*}
%</dis:LaplaceIntro1>
%<*dis:LaplaceIntro2>
注意，第一组方括号内是
下面这个行列式
\begin{equation*}
  \begin{vmat}
    t_{2,2}  &t_{2,3} \\
    t_{3,2}  &t_{3,3}
  \end{vmat}
  \tag{$**$}
\end{equation*}
类似地，其余方括号内也是另外一些行列式。
\begin{equation*}
  \begin{vmat}
    t_{2,1}  &t_{2,3} \\
    t_{3,1}  &t_{3,3}
  \end{vmat}
  \qquad
  \begin{vmat}
    t_{2,1}  &t_{2,2} \\
    t_{3,1}  &t_{3,2}
  \end{vmat}
  \tag{$*\!*\!*$}
\end{equation*}
%</dis:LaplaceIntro2>
%<*dis:LaplaceIntro3>
这些行列式的解释要从置换中寻找。
如果我们固定某一行与某一列\Dash 例如当我们考虑所有含
~$t_{1,1}$ 的项、或所有含~$t_{1,2}$ 的项、或含~$t_{1,3}$ 的项时\Dash 那么 
剩下的就是由置换从其余的
行与列中各选取一个元素。

下面的矩阵对此作了说明。
在左边，含~$t_{1,1}$ 的行与列被涂上阴影，
剩下的正是上面~($**$) 中的矩阵。
类似地，另外两个中剩下的分别是~($*\!*\!*$) 中的矩阵。
\definecolor{shadematrix}{gray}{0.8}
\begin{equation*}
  \left(\begin{array}{>{\columncolor{shadematrix}}ccc}
    \rowcolor{shadematrix}
    t_{1,1}  &t_{1,2}  &t_{1,3} \\
    t_{2,1}  &t_{2,2}  &t_{2,3} \\
    t_{3,1}  &t_{3,2}  &t_{3,3} 
  \end{array}\right)
  \quad
  \left(\begin{array}{c>{\columncolor{shadematrix}}cc}
    \rowcolor{shadematrix}
    t_{1,1}  &t_{1,2}  &t_{1,3} \\
    t_{2,1}  &t_{2,2}  &t_{2,3} \\
    t_{3,1}  &t_{3,2}  &t_{3,3} 
  \end{array}\right)
  \quad
  \left(\begin{array}{cc>{\columncolor{shadematrix}}c}
    \rowcolor{shadematrix}
    t_{1,1}  &t_{1,2}  &t_{1,3} \\
    t_{2,1}  &t_{2,2}  &t_{2,3} \\
    t_{3,1}  &t_{3,2}  &t_{3,3} 
  \end{array}\right)
\end{equation*}
%</dis:LaplaceIntro3>

%<*dis:LaplaceIntro4>
这样，式~($*$) 中出现的那些较小尺寸矩阵的
行列式就有了解释。
至于该式第二行的那个负号，
考虑
$\deter{T}=t_{1,1}t_{2,2}t_{3,3}\deter{P_{\phi_1}}+\cdots+t_{1,3}t_{2,2}t_{3,1}\deter{P_{\phi_6}}$ 的这一写法。
\begin{equation*}
  \begin{aligned}[t]
      &t_{1,1}\cdot \bigg[t_{2,2}t_{3,3}\begin{vmat}[r]
                                 1  &0  &0  \\
                                 0  &1  &0  \\
                                 0  &0  &1
                                \end{vmat}
                 +t_{2,3}t_{3,2}\begin{vmat}[r]
                                  1  &0  &0  \\
                                  0  &0  &1  \\
                                  0  &1  &0
                                \end{vmat}\,\bigg]    \\
         &\hbox{}\quad\hbox{}
          +t_{1,2}\cdot \bigg[t_{2,1}t_{3,3}\begin{vmat}[r]
                                        0  &1  &0  \\
                                        1  &0  &0  \\
                                        0  &0  &1
                                       \end{vmat}
                        +t_{2,3}t_{3,1}\begin{vmat}[r]
                                         0  &1  &0  \\
                                         0  &0  &1  \\
                                         1  &0  &0
                                       \end{vmat}\,\bigg]  \\
        &\hbox{}\quad\hbox{}
         +t_{1,3}\cdot \bigg[t_{2,1}t_{3,2}\begin{vmat}[r]
                                        0  &0  &1  \\
                                        1  &0  &0  \\
                                        0  &1  &0
                                       \end{vmat}
                        +t_{2,2}t_{3,1}\begin{vmat}[r]
                                         0  &0  &1  \\
                                         0  &1  &0  \\
                                         1  &0  &0
                                       \end{vmat}\,\bigg]
    \end{aligned}
\end{equation*}
第一行中的矩阵 $P_{\phi_1}$ 与~$P_{\phi_2}$
恰好是为行列式~($**$) 配置的，
因为它们的第一行与单位矩阵相同，而其余两行
看起来正像该行列式所对应的那两行。
第二行中的矩阵 $P_{\phi_3}$ 与~$P_{\phi_4}$
也可以用同样的方式配置，只需一次
行对换，把~$\rowvec{1 &0 &0}$ 换到第一行。
当然，这次行对换给那一行带来了一个负号。
最后，按同样方式配置第三行的矩阵
需要两次
对换，于是符号保持为正。
%</dis:LaplaceIntro4>

简言之，我们有下面的式子。
\begin{equation*}
  \deter{T}
  =t_{1,1}\cdot \begin{vmat}
            t_{2,2}  &t_{2,3}  \\
            t_{3,2}  &t_{3,3}
          \end{vmat}
   -t_{1,2}\cdot \begin{vmat}
             t_{2,1}  &t_{2,3}  \\
             t_{3,1}  &t_{3,3}
           \end{vmat}
   +t_{1,3}\cdot \begin{vmat}
             t_{2,1}  &t_{2,2}  \\
             t_{3,1}  &t_{3,2}
           \end{vmat}
\end{equation*}
\nearbytheorem{th:LaPlaceExp} 中给出的公式
是上式的推广，
它是一个\definend{递推关系}\index{recurrence}，因为
行列式被表示成一些行列式的组合。
这个公式并非循环论证，因为它把 $\nbyn{n}$
行列式用较小尺寸的矩阵表示出来。

\begin{definition} \label{df:Minor}
%<*df:Minor>
对任意 \( \nbyn{n} \) 矩阵 \( T \)，由删去
\( T \) 的第~\( i \) 行与第~\( j \) 列所得的 \( \nbyn{(n-1)} \) 矩阵是 \( T \) 的
\( i,j \) \definend{子式}\index{minor, of a matrix}\index{determinant!minor}%
\index{matrix!minor} 
。
\( T \) 的 \( i,j \) \definend{代数余子式}\index{cofactor}\index{determinant!using cofactors} 
% \index{matrix!cofactor} 
\( T_{i,j} \) 为
\( (-1)^{i+j} \) 乘以 \( T \) 的 \( i,j \) 子式的行列式。
%</df:Minor>
\end{definition}

\begin{example}
式~($*$) 中那个矩阵的 $1,2$ 代数余子式
是~($*\!*\!*$) 中第一个行列式的相反数。
\begin{equation*}
  T_{1,2}=
  -1\cdot\begin{vmat}
    t_{2,1}  &t_{2,3}  \\
    t_{3,1}  &t_{3,3}
  \end{vmat}
\end{equation*}
\end{example}

\begin{example}
设
\begin{equation*}
   T=
   \begin{mat}[r]
      1  &2  &3  \\
      4  &5  &6  \\
      7  &8  &9
   \end{mat}
\end{equation*}
下面是它的 \( 1,2 \) 与 \( 2,2 \) 代数余子式。
\begin{equation*}
   T_{1,2}=
   (-1)^{1+2}\cdot\begin{vmat}[r]
                4  &6  \\
                7  &9
             \end{vmat}=6
  \qquad
   T_{2,2}=
   (-1)^{2+2}\cdot\begin{vmat}[r]
                1  &3  \\
                7  &9
             \end{vmat}=-12
\end{equation*}
\end{example}

\begin{theorem}[行列式的 Laplace 展开]
\label{th:LaPlaceExp}
\index{determinant!Laplace expansion}%
% \index{Laplace expansion!computes determinant}
%<*th:LaPlaceExp>
设 \( T \) 是 \( \nbyn{n} \) 矩阵，我们可以
通过在任意一行~$i$ 或一列~$j$ 上按代数余子式展开
来求行列式。
\begin{align*}
   \deter{T}
   &=t_{i,1}\cdot T_{i,1}+t_{i,2}\cdot T_{i,2}+\cdots+t_{i,n}\cdot T_{i,n}  \\
   &=t_{1,j}\cdot T_{1,j}+t_{2,j}\cdot T_{2,j}+\cdots+t_{n,j}\cdot T_{n,j}
\end{align*}
%</th:LaPlaceExp>
\end{theorem}

\begin{proof}
证明见 \nearbyexercise{exer:LaplaceProof}。
\end{proof}

\begin{example} \label{ex:ExpLaPlace}
我们可以计算行列式
\begin{equation*}
   \deter{T}=
   \begin{vmat}[r]
     1  &2  &3  \\
     4  &5  &6  \\
     7  &8  &9
   \end{vmat}
\end{equation*}
按第一行展开。
\begin{equation*}
   \deter{T}
   =1\cdot(+1)\begin{vmat}[r]
                5  &6  \\
                8  &9
              \end{vmat}
   +2\cdot(-1)\begin{vmat}[r]
                4  &6  \\
                7  &9
              \end{vmat}
   +3\cdot(+1)\begin{vmat}[r]
                4  &5  \\
                7  &8
              \end{vmat}     
  =-3+12-9     
  =0
\end{equation*}
我们也可以按第二列展开。
\begin{equation*}
   \deter{T}
   =2\cdot(-1)\begin{vmat}[r]
                4  &6  \\
                7  &9
              \end{vmat}
   +5\cdot(+1)\begin{vmat}[r]
                1  &3  \\
                7  &9
              \end{vmat}
   +8\cdot(-1)\begin{vmat}[r]
                1  &3  \\
                4  &6
              \end{vmat}     
  =12-60+48   
  =0
\end{equation*}
\end{example}

\begin{example}
含有很多零的行或列会使这种展开更容易。
\begin{equation*}
  \begin{vmat}[r]
    1 &5  &0  \\
    2 &1  &1  \\
    3 &-1 &0
  \end{vmat}
   =
   0\cdot(+1)\begin{vmat}[r]
               2  &1  \\
               3  &-1
             \end{vmat}+
   1\cdot(-1)\begin{vmat}[r]
               1  &5  \\
               3  &-1
             \end{vmat}+
   0\cdot(+1)\begin{vmat}[r]
               1  &5  \\
               2  &1
             \end{vmat}
  =16
\end{equation*}
\end{example}

最后，我们应用 Laplace 展开
来推导矩阵逆的一个新公式。
利用 \nearbytheorem{th:LaPlaceExp}，我们可以计算 
一个矩阵的行列式， 
方法是取来自某一行的元素与
其相应代数余子式的线性组合。
\begin{equation*}
  t_{i,1}\cdot T_{i,1}+t_{i,2}\cdot T_{i,2}+\dots+t_{i,n}\cdot T_{i,n}
   =\deter{T}  
\end{equation*}
回忆一下，有两行相同的矩阵其行列式为零。
于是，
若用来自
第~$k$ 行（$k\neq i$）的元素为代数余子式加权则得到零
\begin{equation*}
  t_{i,1}\cdot T_{k,1}+t_{i,2}\cdot T_{k,2}+\dots+t_{i,n}\cdot T_{k,n}=0
\end{equation*}
因为它表示的是某个矩阵按第~$k$ 行的展开，而该矩阵的 
第~\( i \) 行等于第~\( k \) 行。
总结如下。
\begin{equation*}
 \generalmatrix{t}{n}{n}
 \begin{mat}
   T_{1,1}  &T_{2,1}  &\ldots  &T_{n,1}  \\
   T_{1,2}  &T_{2,2}  &\ldots  &T_{n,2}  \\
            &\vdots   &        &         \\
   T_{1,n}  &T_{2,n}  &\ldots  &T_{n,n}
 \end{mat}                                  
 =\begin{mat}
     |T|      &0        &\ldots  &0        \\
     0        &|T|      &\ldots  &0        \\
              &\vdots   &        &         \\
     0        &0        &\ldots  &|T|
   \end{mat} 
\end{equation*}
注意，代数余子式矩阵中下标的顺序
与另一个矩阵中下标的顺序相反；例如，
在代数余子式矩阵的第一行中 
下标依次是 $1,1$、$2,1$，等等。

\begin{definition}  \label{df:Adjoint}
%<*df:Adjoint>
方阵 \( T \) 的 \definend{伴随}\index{matrix!adjoint}\index{adjoint matrix}
（或 \definend{伴随矩阵}\index{matrix!adjugate}\index{adjugate matrix}，
或称
\definend{经典伴随}）\index{classical adjoint} 
是
\begin{equation*}
  \adj(T)=
    \begin{mat}
      T_{1,1}  &T_{2,1}  &\ldots  &T_{n,1}  \\
      T_{1,2}  &T_{2,2}  &\ldots  &T_{n,2}  \\
               &\vdots   &        &         \\
      T_{1,n}  &T_{2,n}  &\ldots  &T_{n,n}
    \end{mat}
\end{equation*}
其中第~\( i \) 行第~\( j \) 列处的元素 \( T_{j,i} \)， 
是 \( j,i \) 代数余子式。
%</df:Adjoint>
\end{definition}

\begin{theorem} \label{th:MatTimesAdjEqDiagDets}
%<*th:MatTimesAdjEqDiagDets>
设 \( T \) 是方阵，
$T\cdot \adj(T)=\adj(T)\cdot T=\deter{T}\cdot I$。
于是，若 $T$ 有逆，即 $\deter{T}\neq 0$，则
$T^{-1}=(1/\deter{T})\cdot\adj(T)$。\index{matrix!inverse}\index{inverse!matrix}
%</th:MatTimesAdjEqDiagDets>
\end{theorem}

\begin{proof}
定理表述之前的讨论已经把这一点讲清楚了。
\end{proof}

\begin{example}    \label{ex:MatTimesAdjEqDets}
若
\begin{equation*}
  T=\begin{mat}[r]
        1  &0  &4  \\
        2  &1  &-1 \\
        1  &0  &1
      \end{mat}
\end{equation*}
则 $\adj(T)$ 为
\begin{equation*}
    \begin{mat}
      T_{1,1}  &T_{2,1}  &T_{3,1} \\
      T_{1,2}  &T_{2,2}  &T_{3,2} \\
      T_{1,3}  &T_{2,3}  &T_{3,3}   
    \end{mat}
    \!\!=\!\!\begin{mat}
      \begin{vmat}[r]
           1  &-1 \\
           0  &1
         \end{vmat}
      &-\begin{vmat}[r]
             0  &4  \\
             0  &1
           \end{vmat}
      &\begin{vmat}[r]
             0  &4  \\
             1  &-1
         \end{vmat}             \\[2.1ex]
      -\begin{vmat}[r]
           2  &-1 \\
           1  &1
         \end{vmat}
      &\begin{vmat}[r]
             1  &4  \\
             1  &1
           \end{vmat}
      &-\begin{vmat}[r]
             1  &4  \\
             2  &-1
           \end{vmat}            \\[2.1ex]
      \begin{vmat}[r]
           2  &1  \\
           1  &0
         \end{vmat}
      &-\begin{vmat}[r]
             1  &0  \\
             1  &0
           \end{vmat}
      &\begin{vmat}[r]
             1  &0  \\
             2  &1
           \end{vmat}
    \end{mat}
    \!\!=\!                 
    \begin{mat}[r]
       1  &0  &-4  \\
       -3 &-3 &9  \\
       -1 &0  &1
    \end{mat}         
\end{equation*}
而它与 $T$ 作乘积则给出对角矩阵 $\deter{T}\cdot I$。
\begin{equation*}
  \begin{mat}[r]
    1  &0  &4  \\
    2  &1  &-1 \\
    1  &0  &1
  \end{mat}
  \begin{mat}[r]
    1  &0  &-4  \\
    -3 &-3 &9  \\
    -1 &0  &1
   \end{mat}         
  =\begin{mat}[r]
     -3  &0  &0  \\
      0  &-3 &0  \\
      0  &0  &-3
   \end{mat}
\end{equation*}
% \end{example}

% \begin{corollary} \label{cor:InvFromAdj}
% \index{matrix!inverse}\index{inverse!matrix}
% %<*co:InvFromAdj>
% If \( \deter{T}\neq 0 \) then
% $T^{-1}=(1/\deter{T})\cdot\adj(T)$.
% %</co:InvFromAdj>
% \end{corollary}

% \begin{example}
\noindent $T$ 的逆 
% The inverse of the matrix from \nearbyexample{ex:MatTimesAdjEqDets}
是 $(1/-3)\cdot\adj(T)$。
\begin{equation*}
  T^{-1}
  =\begin{mat}[r]  % braces make the - a negative sign?
     1/\hbox{$-3$}  &0/\hbox{$-3$}  &-4/\hbox{$-3$}  \\
    -3/\hbox{$-3$}  &-3/\hbox{$-3$} &9/\hbox{$-3$}   \\
    -1/\hbox{$-3$}  &0/\hbox{$-3$}  &1/\hbox{$-3$}
   \end{mat}
  =\begin{mat}[r]
     -1/3  &0  &4/3  \\
      1    &1  &-3   \\
      1/3  &0  &-1/3
   \end{mat}
\end{equation*}
\end{example}

本小节的公式常用于手工
计算，并且有时对特殊类型的矩阵也很有用。
然而，对于一般的矩阵，它们
并不是最好的选择， 
因为与例如 
Gauss–Jordan 方法相比，它们需要更多的算术运算。




